\chapter{Nociones de lógica}
\label{cap:logica}

\section{Proposiciones y conectivas}

La \textbf{sintaxis} estudia cómo se forman las expresiones; la
\textbf{semántica}, su significado y valor de verdad. Una
\textbf{proposición} es un enunciado al que se puede asignar un valor
de verdad: verdadero ($1$) o falso ($0$). Puede ser simple o estar
compuesta por otras proposiciones.

Un \textbf{predicado} contiene variables y se convierte en proposición
al fijar sus valores en un dominio. Para $x$ entero, «$x$ es múltiplo de $5$»
depende de $x$; para $x=15$ es verdadero y para $x=12$, falso.
Una relación puede depender de varias variables.

Las \textbf{conectivas} construyen proposiciones compuestas. Su valor
depende únicamente de los valores de sus argumentos:
\begin{itemize}
 \item $\neg p$: negación, que invierte el valor de $p$.
 \item $p\vee q$: disyunción inclusiva, verdadera si al menos una es verdadera.
 \item $p\wedge q$: conjunción, verdadera si ambas son verdaderas.
 \item $p\to q$: condicional, falso únicamente cuando $p$ es verdadero y
       $q$ es falso. $p$ es el antecedente y $q$, el consecuente.
 \item $p\leftrightarrow q$: bicondicional, verdadero si ambos valores coinciden.
\end{itemize}
No se exige una relación causal para evaluar un condicional. Si su
antecedente es falso, el condicional es verdadero.

\begin{center}
\begin{tabular}{@{}ccccccc@{}}
\toprule
$p$ & $q$ & $\neg p$ & $p\vee q$ & $p\wedge q$ & $p\to q$ & $p\leftrightarrow q$\\
\midrule
1&1&0&1&1&1&1\\
1&0&0&1&0&0&0\\
0&1&1&1&0&1&0\\
0&0&1&0&0&1&1\\
\bottomrule
\end{tabular}
\end{center}

Una tabla de verdad de $n$ variables tiene $2^n$ filas. Hay
$2^{2^n}$ funciones posibles de esos valores: cuatro conectivas unarias
y dieciséis binarias. En el caso unario son $0$, $1$, $p$ y $\neg p$.
La disyunción exclusiva es verdadera cuando los dos valores son distintos.

Los paréntesis indican qué argumentos recibe cada conectiva.
No deben omitirse si ello permite interpretaciones diferentes.

Para $p\to q$ se distinguen la \textbf{recíproca} $q\to p$,
la \textbf{contraria} $\neg p\to\neg q$ y la
\textbf{contrarrecíproca} $\neg q\to\neg p$.
Solo esta última es siempre equivalente al condicional original.
La recíproca y la contraria son equivalentes entre sí.

\section{Equivalencias y simplificación}
\label{sec:equivalencias}

Una \textbf{tautología} es verdadera en todas las asignaciones;
una \textbf{contradicción}, falsa en todas. Se representan por $1$ y $0$,
respectivamente. Las demás fórmulas cambian de valor según la asignación.

Dos fórmulas $P$ y $Q$ son \textbf{equivalentes}, $P\eqv Q$, si coinciden
en todas las asignaciones, es decir, si $P\leftrightarrow Q$ es una tautología.
Coincidir en una sola fila de la tabla no basta.
La notación $P\imp Q$ expresa una implicación válida en todas las
asignaciones: $P\to Q$ es una tautología. Los signos $\to$ y
$\leftrightarrow$ forman fórmulas; $\imp$ y $\eqv$ afirman su validez.

Se puede sustituir una subfórmula por otra equivalente sin cambiar el
valor de la expresión completa. Las siguientes leyes permiten simplificar:
\begin{align*}
 \neg\neg p&\eqv p, &
 p\vee p&\eqv p, & p\wedge p&\eqv p,\\
 p\vee0&\eqv p, &
 p\wedge1&\eqv p, & p\vee1&\eqv1,\\
 p\wedge0&\eqv0, &
 p\vee\neg p&\eqv1, & p\wedge\neg p&\eqv0.
\end{align*}
La disyunción y la conjunción son conmutativas y asociativas.
El bicondicional también lo es; el condicional no permite intercambiar
ni reagrupar sus argumentos en general.

Las leyes distributivas conectan conjunción y disyunción:
\begin{align*}
 p\wedge(q\vee r)&\eqv(p\wedge q)\vee(p\wedge r),\\
 p\vee(q\wedge r)&\eqv(p\vee q)\wedge(p\vee r).
\end{align*}
Las \textbf{leyes de De Morgan} indican cómo negar cada una:
\[
 \neg(p\wedge q)\eqv\neg p\vee\neg q,\qquad
 \neg(p\vee q)\eqv\neg p\wedge\neg q.
\]
El condicional y el bicondicional se pueden expresar con otras conectivas:
\begin{align*}
 p\to q&\eqv\neg p\vee q\eqv\neg(p\wedge\neg q),\\
 p\leftrightarrow q&\eqv(p\to q)\wedge(q\to p)
                 \eqv(p\wedge q)\vee(\neg p\wedge\neg q).
\end{align*}
Además,
\[
 p\to q\eqv\neg q\to\neg p,\qquad
 p\leftrightarrow q\eqv\neg p\leftrightarrow\neg q.
\]
Estas equivalencias explican el uso de la contrarrecíproca.
Para una contradicción $q\wedge\neg q$,
\[
 \neg p\to(q\wedge\neg q)\eqv p,
\]
que es la base lógica de la reducción al absurdo.

\section{Inferencias y comprobación de argumentos}
\label{sec:inferencias}

Un argumento es \textbf{válido} si no existe una asignación que haga
verdaderas todas las premisas y falsa la conclusión. Con premisas
$P_1,\ldots,P_k$ y conclusión $Q$, se comprueba la tautología
\[
 (P_1\wedge\cdots\wedge P_k)\to Q.
\]
La validez conserva la verdad de las premisas; no afirma por sí sola
que estas sean verdaderas.

\begin{center}
\small
\begin{tabular}{@{}p{0.31\linewidth}p{0.38\linewidth}p{0.20\linewidth}@{}}
\toprule
Regla & Premisas & Conclusión\\
\midrule
Simplificación & $p\wedge q$ & $p$ (o $q$)\\
Ampliación & $p$ & $p\vee q$\\
Modus ponens & $p$, $p\to q$ & $q$\\
Modus tollens & $\neg q$, $p\to q$ & $\neg p$\\
Silogismo disyuntivo & $\neg p$, $p\vee q$ & $q$\\
Silogismo disyuntivo dual & $p$, $\neg p\vee\neg q$ & $\neg q$\\
\bottomrule
\end{tabular}
\end{center}

Para verificar una equivalencia o una inferencia se puede:
\begin{itemize}
 \item Construir la tabla completa de verdad y comparar las columnas
       o comprobar que el condicional correspondiente siempre vale $1$.
 \item Transformar las fórmulas mediante las leyes de
       \S\ref{sec:equivalencias}, o encadenar inferencias válidas.
 \item Intentar una \textbf{refutación}: suponer verdaderas las premisas
       y falsa la conclusión. Si esos valores fuerzan una contradicción,
       el argumento es válido; si son compatibles, constituyen un contraejemplo.
\end{itemize}

\section{Forma normal conjuntiva}
\label{sec:fnc}

Un \textbf{literal} es una variable proposicional o su negación.
Una \textbf{cláusula} es una disyunción de literales.
Una fórmula está en \textbf{forma normal conjuntiva} (FNC) cuando es una
conjunción de cláusulas. Para que sea verdadera deben ser verdaderas
todas las cláusulas; para que una cláusula lo sea basta un literal verdadero.

Para obtener una FNC equivalente:
\begin{enumerate}
 \item Elimina los bicondicionales y después los condicionales usando
       las equivalencias de \S\ref{sec:equivalencias}.
 \item Aplica De Morgan y doble negación hasta que cada negación
       afecte solo a una variable.
 \item Distribuye $\vee$ sobre $\wedge$ para obtener una conjunción
       de disyunciones de literales.
 \item Simplifica mediante las leyes lógicas cuando sea posible.
\end{enumerate}

\begin{ejemplo}
La fórmula $\neg(p\to q)\vee r$ se transforma así:
\[
 \neg(p\to q)\vee r
 \eqv(p\wedge\neg q)\vee r
 \eqv(p\vee r)\wedge(\neg q\vee r).
\]
La última expresión tiene dos cláusulas y está en FNC.
\end{ejemplo}

La FNC no es única: dos formas diferentes pueden representar la misma
función de verdad. Obtener la misma FNC sí garantiza equivalencia.

\section{Axiomas y demostraciones}

\subsection*{Sistema PM}

Un sistema axiomático fija símbolos, axiomas y reglas de deducción.
En PM las conectivas primitivas son $\neg$ y $\vee$; las demás se definen
a partir de ellas:
\[
 p\wedge q:=\neg(\neg p\vee\neg q),\qquad p\to q:=\neg p\vee q,
\]
\[
 p\leftrightarrow q:=(p\to q)\wedge(q\to p).
\]
Sus cuatro esquemas de axiomas son:
\begin{align*}
 &(p\vee p)\to p,\\
 &p\to(p\vee q),\\
 &(p\vee q)\to(q\vee p),\\
 &(p\to q)\to\bigl[(r\vee p)\to(r\vee q)\bigr].
\end{align*}
La \textbf{sustitución} reemplaza una variable por una fórmula en todas
sus apariciones. La \textbf{separación} aplica modus ponens: de $P$
y $P\to Q$ se deduce $Q$. Una demostración formal es una sucesión
de fórmulas obtenidas de los axiomas mediante estas reglas.

\subsection*{Resultados y condiciones}

Un \textbf{axioma} se acepta como punto de partida; un
\textbf{teorema} es un resultado demostrado. Se llama
\textbf{lema} a un resultado auxiliar, \textbf{corolario} a una consecuencia
directa y \textbf{proposición} a un resultado que no suele destacarse como
teorema. Estas denominaciones organizan la exposición.
Una \textbf{conjetura} es una afirmación propuesta cuya demostración
aún no se conoce.

En una implicación $P\imp Q$, $P$ es condición \textbf{suficiente} para
$Q$ y $Q$ es condición \textbf{necesaria} para $P$.
Si $P\eqv Q$, cada una es necesaria y suficiente para la otra:
la equivalencia da una caracterización.

\subsection*{Métodos de demostración}

\begin{itemize}
 \item \textbf{Directa:} parte de las hipótesis y encadena consecuencias
       justificadas hasta la conclusión.
 \item \textbf{Por contrarrecíproca:} para demostrar $P\imp Q$,
       demuestra $\neg Q\imp\neg P$.
 \item \textbf{Por reducción al absurdo:} para demostrar $Q$ bajo unas
       hipótesis, añade $\neg Q$ y deduce una contradicción.
\end{itemize}
En todos los casos hay que distinguir lo que se supone de lo que se
pretende demostrar y justificar cada paso con definiciones o resultados válidos.
